\section{Related Work}
\label{sec:related-work}

Our work brings together end-user programming of AI behavior, interactive
control of agent workflows, and maintenance of LLM-based systems. These areas
describe how people express behavior through reusable natural-language
artifacts, inspect the resulting workflows, and revise them after breakdowns.

\subsection{End-User Programming of AI Behavior}
\label{sec:rw-eup}

End-user programming enables domain experts to shape computational behavior
through task-level representations without mastering conventional programming
\cite{ko2011state}. Reusable natural-language instructions extend this premise
to LLMs: they are inspectable and editable, yet their operational effects depend
on model interpretation, surrounding context, available tools, and sampling.
Prompt-authoring systems such as ChainForge and EvalLM therefore supplement
text editing with comparative generation and user-defined evaluation criteria
\cite{arawjo2024chainforge,kim2024evallm}. Policy-authoring systems broaden the
unit of design from individual outputs to regions of behavior. Policy Maps
supports exploration of an open-ended behavior space, while Data-Prompt
Co-Evolution couples prompt revision with edge cases, rationales, neighboring
examples, and a living test set \cite{lam2025policymaps,lee2026dataprompt}.
These systems show how examples can articulate behavioral boundaries and
evaluate an existing policy. This framing is particularly relevant when the
person accountable for a workflow has domain authority without maintaining
every underlying instruction. A readable natural-language artifact can support
control while still leaving its runtime behavior uncertain.

We study these natural-language artifacts after they have entered repeated use.
A breakdown motivates a \emph{relative} change to a skill that already produces
behavior its owner values. The owner must decide what should improve, what
should remain stable, and which instructions implement that boundary.
Additional cases may surface tacit commitments or unresolved trade-offs,
allowing the repair specification to evolve with the evidence.

\subsection{Interactive Control of Agent Workflows}
\label{sec:rw-interactive-control}

Interactive agent systems make runtime activity more actionable through shared
representations and counterfactual intervention. Cocoa uses an interactive plan
to support co-planning, delegation, execution, and replanning between a user and
an agent \cite{feng2026cocoa}. Agent-debugging systems expose complementary
views of execution: AGDebugger supports navigation, editing, and reset of
multi-agent conversation branches, while RAG Without the Lag enables rapid
what-if analysis over retrieval and generation components
\cite{epperson2025agdebugger,romerolauro2026rag}. These systems show how
representations connected to runtime actions and alternatives can make agent
activity more inspectable and steerable than an undifferentiated chat history.
They also establish a useful interaction pattern: users can intervene in a
concrete execution without reconstructing the agent's full internal process.

These techniques primarily support the coordination, interpretation, or
steering of a current task. A trace or branch comparison offers evidence about
one enacted workflow. Persistent repair also requires evidence about stability
across runs, transfer across contexts, and alignment with the intended scope.
Research on appropriate reliance further shows that sources and visible
inconsistencies can help users recognize uncertainty in AI outputs
\cite{kim2025reliance}. We extend interactive control from one trajectory to a
revision that may shape a family of future workflows.

\subsection{Testing and Maintenance of LLM-Based Systems}
\label{sec:rw-maintenance}

Software maintenance distinguishes repairing a motivating failure from
protecting behavior that previously worked. Regression testing checks prior
behavior after a change, while code review makes a proposed patch and its
rationale inspectable before integration
\cite{rothermel1996regression,bacchelli2013review}. These practices become less
direct for LLM artifacts: executions are probabilistic and context-dependent,
and textual locality does not guarantee behavioral locality. PromptPex extracts
input--output specifications from a prompt and generates tests that can expose
invalid behavior and later regressions \cite{sharma2025promptpex}. Contrastive
Reflection combines failures with nearby successes to propose targeted edits and
select candidates using held-out validation and optional regression constraints
\cite{koh2026contrastive}. These methods show how prompt repair can use tests
and regression constraints to evaluate more than a satisfactory rerun.
PromptPex is most decisive when the prompt already states the relevant contract,
while Contrastive Reflection uses held-out examples or task-level criteria to
rank candidate edits. When the artifact leaves the policy incomplete,
conflicting, or overly broad, additional evidence may reshape the repair scope
as well as evaluate a candidate.

Recent systems extend execution-grounded repair from prompts to persistent
skills and agent harnesses. SkillRevise uses execution traces, candidate
revisions, re-execution, and empirical utility to improve an initial skill;
SkillHEX turns falsifiable failure hypotheses into executable tests and searches
persistent revision branches \cite{liu2026skillrevise,feng2026skillhex}.
HarnessFix localizes failed trajectories across layers of an agent harness,
applies scoped repair operators, and validates patches on held-out tasks and
regressions \cite{chen2026harnessfix}. These systems demonstrate
execution-grounded diagnosis, revision search, and regression validation across
several types of agent artifacts. We focus on the reusable skill as the
owner-facing repair object. Empirical utility, executable tests, and held-out
tasks can compare revisions once the target is sufficiently specified.
Workflow owners may still be forming the intended policy and deciding which
instructions should implement it.

Skill repair often begins while the desired policy remains unsettled. A
workflow owner may reject one outcome while still deciding how neighboring
cases should behave, which instructions deserve attention, and what evidence
would justify reuse. Additional executions can refine the specification or
reveal effects beyond the motivating case. We study owner-governed skill
maintenance, complementing automatic repair with support for forming a relative
behavioral specification, connecting instructions to cross-context evidence,
and deciding whether to release a persistent change.
